\documentclass[a4paper,11pt]{article}
\usepackage[T2A]{fontenc}
\usepackage[utf8]{inputenc}
\usepackage[russian]{babel}
\usepackage{amsmath,amssymb,mathtools}
\usepackage{helvet}
\renewcommand{\familydefault}{\sfdefault}
\usepackage[margin=20mm,top=20mm,bottom=22mm]{geometry}
\usepackage{microtype}
\usepackage{booktabs,tabularx,array}
\usepackage{enumitem}
\usepackage{xcolor}
\usepackage{hyperref}
\usepackage{fancyhdr}
\usepackage{tikz}
\usepackage{icomma}
\definecolor{accent}{HTML}{086A74}
\definecolor{muted}{HTML}{515D62}
\hypersetup{colorlinks=true,linkcolor=accent,urlcolor=accent,pdfauthor={Лёвин Артём Александрович},pdftitle={Ревью домашней работы — Анастасия Павлова, 26 октября 2026}}
\setlength{\parindent}{0pt}
\setlength{\parskip}{5pt}
\setlength{\headheight}{14pt}
\raggedbottom
\pagestyle{fancy}
\fancyhf{}
\fancyhead[L]{\footnotesize\color{muted} Математика · Банковские вклады и кредиты}
\fancyhead[R]{\footnotesize\color{muted}26.10.2026}
\fancyfoot[C]{\scriptsize Лёвин Артём Александрович · эксклюзивно для Анастасии Павловой\quad | \thepage}
\renewcommand{\headrulewidth}{0.3pt}
\newcommand{\err}{\textcolor{red!75!black}{\textbf{Ошибка.}}\ }
\newcommand{\idea}{\textcolor{accent}{\textbf{Ключевая идея.}}\ }
\newcommand{\result}{\textcolor{accent}{\textbf{Итог.}}\ }
\newcommand{\task}[2]{\clearpage\phantomsection\label{task:#1}\section*{Задание №\,#1\;·\;#2}}
\newcommand{\labelp}[1]{\par\smallskip\textbf{#1.}\enspace}
\begin{document}
\thispagestyle{empty}
\vspace*{35mm}
\begin{center}
{\large\color{muted} МАТЕМАТИКА}\par\vspace{12mm}
{\Huge\bfseries Ревью домашней работы}\par\vspace{7mm}
{\Large Банковские вклады и кредиты}\par\vspace{20mm}
{\Large Анастасия Павлова}\par\vspace{13mm}
{\large 26 октября 2026 года}\par
\end{center}
\vfill
\begin{center}
\textbf{Лёвин Артём Александрович}\par
эксклюзивно для Анастасии Павловой
\end{center}
\vspace{10mm}
\begin{center}
\begin{tikzpicture}
\draw[accent,opacity=.65,line width=1pt] (-5,0) .. controls (-2.5,.60) and (-.8,-.48) .. (1,0) .. controls (2.2,.34) and (3.7,-.18) .. (5,0);
\end{tikzpicture}
\end{center}
\clearpage
\section*{Сводная таблица ошибок}
В таблице представлены только задания, требующие исправления. Номер задания открывает подробный разбор.

\renewcommand{\arraystretch}{1.3}
\small
\begin{tabularx}{\linewidth}{@{}p{11mm}p{23mm}X p{24mm}p{28mm}@{}}
\toprule
\textbf{№} & \textbf{Тема} & \textbf{Суть ошибки} & \textbf{Ваш ответ} & \textbf{Верный ответ}\\\midrule
\hyperref[task:2]{2} & Вклад, ставка & Ошибка при раскрытии куба суммы & Записан неверный способ; ставка не найдена & $10\%$\\
\hyperref[task:4]{4} & Пополнения & Ошибка сложения начислений & 166\,200 руб. & 156\,200 руб.\\
\hyperref[task:5]{5} & Сравнение вкладов & Неверное сравнение итоговых сумм и единиц & Б выгоднее на 19 тыс. руб. & А выгоднее на 990 руб.\\
\hyperref[task:6]{6} & Два платежа & Неверный исходный кредит; пропущен платёж & $\approx59\,182$ руб. & 121\,000 руб.\\
\hyperref[task:7]{7} & Три платежа & Не учтён третий платёж & $\approx286\,078$ руб. & 199\,650 руб.\\
\hyperref[task:8]{8} & Сумма кредита & Не учтён последний платёж & 121\,000 руб. & 331\,000 руб.\\
\hyperref[task:9]{9} & Переплата & Опора на неверный платёж, знак и отсутствие процента & $-361\,734$ руб.; процент не указан & 102\,450 руб.; $20{,}63\%$\\
\hyperref[task:10]{10} & Четыре платежа & Решение не представлено & Не указан & 464\,100 руб.\\
\bottomrule
\end{tabularx}
\normalsize
\vfill
\textcolor{muted}{\footnotesize Для задач с ежегодными равными платежами считаем, что в конце каждого года сначала начисляются проценты на остаток долга, затем вносится очередной платёж.}
\clearpage
\thispagestyle{empty}
\vspace*{.37\textheight}
\begin{center}{\Huge\bfseries\color{accent} Разбор ошибок}\end{center}

\task{2}{Определение годовой ставки}
\labelp{Текст задания} На вклад внесли 50\,000 руб. Через три года при ежегодной капитализации получили 66\,550 руб. Найдите постоянную годовую ставку.
\labelp{Ваше решение} Вы верно начинаете с равенства $66{,}55=50(1+r)^3$, где суммы записаны в тысячах рублей. Далее в записях появляется $16{,}55=50r^3$ и $r^3=0{,}331$. Затем ставка пытается вычисляться с помощью кубического корня и преобразования процентов.
\labelp{Ваш ответ} Достоверного окончательного значения ставки в процентах в работе нет.
\labelp{Где именно возникает ошибка} Последний правильный шаг:
\[66{,}55=50(1+r)^3.\]
Первый неправильный переход:
\[66{,}55-50=50r^3.\]
\err При раскрытии $(1+r)^3$ потеряны слагаемые $3r$ и $3r^2$. Верно
\[(1+r)^3=1+3r+3r^2+r^3.\]
Следовательно, из $50(1+r)^3-50=16{,}55$ нельзя заключать, что $50r^3=16{,}55$. Дальнейшее извлечение корня из $0{,}331$ основано на уже нарушенном равенстве.
\labelp{Исправление} Вернитесь к последнему корректному равенству; сначала разделите обе части на 50:
\[(1+r)^3=\frac{66{,}55}{50}=1{,}331.\]
\idea Коэффициент роста за три года равен кубу коэффициента роста за один год. Поэтому извлечь кубический корень нужно из \emph{всего} $1{,}331$.
\labelp{Полное правильное решение}
\[(1+r)^3=1{,}331=1{,}1^3.\]
Поскольку $1+r>0$, получаем $1+r=1{,}1$, откуда $r=0{,}1$. Переводим долю в проценты: $0{,}1\cdot100\%=10\%$.
\labelp{Проверка} $50\,000\cdot1{,}1^3=50\,000\cdot1{,}331=66\,550$ руб. Начисление совпадает с условием.
\result Годовая ставка составляет $10\%$.
\labelp{Задача на закрепление} На вклад внесли 25\,000 руб. Через три года с ежегодной капитализацией на нём оказалось 43\,200 руб. Найдите постоянную годовую процентную ставку.

\task{4}{Вклад с ежегодными пополнениями}
\labelp{Текст задания} На вклад внесли 100\,000 руб. под $10\%$ годовых. После начисления процентов в конце первого и второго года добавили по 10\,000 руб. Сколько будет на вкладе в конце третьего года?
\labelp{Ваше решение} Вы учитываете исходный вклад за три года и два пополнения за оставшиеся два года и один год соответственно:
\[100\cdot1{,}1^3+10\cdot1{,}1^2+10\cdot1{,}1.\]
В тысячах рублей получены слагаемые $133{,}1$, $12{,}1$ и $11$, после чего в записи стоит сумма $166{,}2$.
\labelp{Ваш ответ} 166\,200 руб.
\labelp{Где именно возникает ошибка} Последний правильный шаг:
\[133{,}1+12{,}1+11.\]
Первый неверный шаг --- арифметическое сложение этих трёх величин:
\[133{,}1+12{,}1+11\ne166{,}2.\]
\err При верно составленной формуле допущена арифметическая ошибка. Сначала $133{,}1+12{,}1=145{,}2$, затем $145{,}2+11=156{,}2$.
\idea Каждое пополнение начисляет проценты только за оставшееся до окончания вклада время. Деньги, добавленные после первого года, работают два года; после второго --- один год.
\labelp{Исправление} Сохраняем Вашу формулу и исправляем последнюю операцию:
\[133\,100+12\,100+11\,000=156\,200.\]
\labelp{Полное правильное решение} После первого года: $100\,000\cdot1{,}1+10\,000=120\,000$ руб. После второго года: $120\,000\cdot1{,}1+10\,000=142\,000$ руб. После третьего: $142\,000\cdot1{,}1=156\,200$ руб.
\labelp{Проверка} Независимый подсчёт по трём взносам даёт $100\,000\cdot1{,}331+10\,000\cdot1{,}21+10\,000\cdot1{,}1=156\,200$ руб.
\result В конце третьего года на вкладе 156\,200 руб.
\labelp{Задача на закрепление} На вклад внесли 80\,000 руб. под $10\%$ годовых. В конце первого и второго года после начисления процентов пополнили вклад на 5\,000 руб. Найдите сумму в конце третьего года.

\task{5}{Сравнение двух вариантов вклада}
\labelp{Текст задания} Вклад --- 1\,000 тыс. руб., срок три года. Вариант А: $10\%$ ежегодно. Вариант Б: $7\%$, $10\%$, $13\%$ по годам. Какой выгоднее и на сколько рублей?
\labelp{Ваше решение} В таблице Вы последовательно находите сумму по годам для двух вариантов: для А доходите до $1331$ тыс. руб., для Б --- приблизительно до $1330$ тыс. руб. Затем появляется вычитание $1330-1311=19$ и вывод о преимуществе Б.
\labelp{Ваш ответ} Вариант Б выгоднее на 19 тысяч рублей.
\labelp{Где именно возникает ошибка} Последний верный этап --- расчёт сумм по годам с применением ежегодных коэффициентов. Первый ошибочный переход --- сравнение конечных сумм: вместо уже полученных $1331$ и примерно $1330$ в вычитании использовано число $1311$; кроме того, выбран неверный порядок сравнения.
\err Неверно переписана итоговая сумма варианта А, поэтому изменились и разность, и вывод о выгоде. Для ответа в рублях необходимо также перевести тысячи рублей в рубли.
\idea При меняющихся ставках коэффициенты \emph{перемножают}, потому что каждый следующий процент начисляется на результат предыдущего года.
\labelp{Исправление} Сравните две конечные суммы без промежуточных округлений:
\[S_{\mathrm A}=1\,000\,000\cdot1{,}1^3,\qquad
S_{\text{Б}}=1\,000\,000\cdot1{,}07\cdot1{,}10\cdot1{,}13.\]
\labelp{Полное правильное решение}
\begin{align*}
S_{\mathrm A}&=1\,000\,000\cdot1{,}331=1\,331\,000\text{ руб.},\\
S_{\text{Б}}&=1\,000\,000\cdot1{,}33001=1\,330\,010\text{ руб.}
\end{align*}
Разность равна $1\,331\,000-1\,330\,010=990$ руб.
\labelp{Проверка} Оба результата превышают исходный миллион; первый больше второго ровно на 990 руб. Разница составляет $0{,}99$ тыс. руб.
\result Выгоднее вариант А на 990 руб.
\labelp{Задача на закрепление} Вклад 500\,000 руб. размещён на три года. Вариант А: $10\%$ ежегодно. Вариант Б: по годам $8\%$, $9\%$, $13\%$. Какой вариант выгоднее и на сколько рублей?

\task{6}{Кредит с двумя равными платежами}
\labelp{Текст задания} Взят кредит 210\,000 руб. под $10\%$ годовых на два года. Долг погашается двумя равными ежегодными платежами. Найдите один платёж.
\labelp{Ваше решение} В начале записи используется число $240$ (тыс. руб.) вместо $210$. Далее в уравнении учитывается процентный рост кредита и один вычитаемый платёж, а в ответе указано около 59\,182 руб.
\labelp{Ваш ответ} Приблизительно 59\,182 руб.
\labelp{Где именно возникает ошибка} Правильная идея --- обозначить ежегодный платёж через $x$ и умножать остаток долга на $1{,}1$. Первый достоверный дефект --- замена исходных $210$ тыс. руб. на $240$ тыс. руб. Второй --- отсутствие вычитания платежа в конце второго года.
\err Каждый из двух платежей уменьшает долг. Первый платёж, внесённый в конце первого года, дополнительно уменьшает сумму, на которую начисляются проценты второго года.
\idea Составляйте уравнение по годам: начислили проценты $\rightarrow$ вычли платёж $\rightarrow$ снова начислили проценты $\rightarrow$ вычли второй платёж $\rightarrow$ остаток ноль.
\labelp{Исправление} Обозначим платёж $x$. После первого года остаток $210\,000\cdot1{,}1-x$. После второго года:
\[(210\,000\cdot1{,}1-x)\cdot1{,}1-x=0.\]
\labelp{Полное правильное решение} Раскроем скобки:
\[210\,000\cdot1{,}21-1{,}1x-x=0,\]
\[254\,100-2{,}1x=0,\qquad x=\frac{254\,100}{2{,}1}=121\,000.\]
\labelp{Проверка} Год 1: $210\,000\cdot1{,}1-121\,000=110\,000$. Год 2: $110\,000\cdot1{,}1-121\,000=0$.
\result Каждый ежегодный платёж равен 121\,000 руб.
\labelp{Задача на закрепление} Кредит 315\,000 руб. выдан под $10\%$ годовых на два года с погашением двумя равными платежами в конце каждого года. Найдите платёж.

\task{7}{Кредит с тремя равными платежами}
\labelp{Текст задания} Кредит 496\,500 руб. под $10\%$ годовых взят на три года. Его погашают тремя равными ежегодными платежами. Найдите платёж и общую сумму выплат.
\labelp{Ваше решение} Использован трёхлетний коэффициент $1{,}331$ и коэффициенты $1{,}21$, $1{,}1$ при платежах. Однако в итоговом уравнении Вы собираете при $x$ коэффициент $2{,}31$ и находите приблизительно 286\,078 руб.
\labelp{Ваш ответ} Около 286\,078 руб.; общая сумма платежей отдельно не указана.
\labelp{Где именно возникает ошибка} Последний верный элемент --- коэффициенты $1{,}21$ и $1{,}1$ для двух первых платежей. Первый неверный шаг --- переход к коэффициенту $2{,}31$ вместо $1{,}21+1{,}1+1=3{,}31$.
\err Учтены только первые два платежа. Третий платёж тоже вычитают из долга, хотя после него уже нет начисления процентов. Поэтому он входит в уравнение с коэффициентом $1$.
\idea Если все суммы перенести к концу третьего года, стоимость кредита равна сумме приведённых к этой дате платежей: $P\cdot1{,}1^3=x\cdot1{,}1^2+x\cdot1{,}1+x$.
\labelp{Исправление} Добавьте третий платёж и корректный знак:
\[496\,500\cdot1{,}331=1{,}21x+1{,}1x+x.\]
\labelp{Полное правильное решение}
\[660\,841{,}5=3{,}31x,\qquad x=\frac{660\,841{,}5}{3{,}31}=199\,650.\]
Всего сделаны три платежа, поэтому $3x=3\cdot199\,650=598\,950$ руб.
\labelp{Проверка} После первого года: $546\,150-199\,650=346\,500$ руб. После второго: $381\,150-199\,650=181\,500$ руб. После третьего: $199\,650-199\,650=0$.
\result Платёж --- 199\,650 руб.; общая сумма --- 598\,950 руб.
\labelp{Задача на закрепление} Кредит 662\,000 руб. взят под $10\%$ годовых на три года с тремя одинаковыми ежегодными платежами. Найдите платёж и сумму всех платежей.

\task{8}{Определение первоначальной суммы кредита}
\labelp{Текст задания} Кредит полностью погашен тремя равными ежегодными платежами по 133\,100 руб. Годовая ставка $10\%$. Найдите первоначальную сумму кредита.
\labelp{Ваше решение} Для первоначальной суммы $x$ Вы записываете выражение с $1{,}331x$ и вычитанием выплат с множителями $1{,}21$ и $1{,}1$. В последующих преобразованиях появляется $1{,}331x=161\,051$ и итог 121\,000 руб.
\labelp{Ваш ответ} 121\,000 руб.
\labelp{Где именно возникает ошибка} Верно найден коэффициент роста кредита за три года: $1{,}331$. Первый дефект модели --- нет вычитания \emph{третьего} платежа $133\,100$ руб. В дальнейших вычислениях также теряется вклад второго платежа.
\err При трёх платежах в уравнении должны присутствовать все три слагаемых: первый платёж накапливается ещё два года, второй --- один, третий берётся без дополнительного роста.
\idea Можно составить уравнение из будущих стоимостей всех трёх платежей на дату окончания кредита.
\labelp{Исправление}
\[1{,}331x=133\,100\cdot1{,}21+133\,100\cdot1{,}1+133\,100.\]
\labelp{Полное правильное решение}
\begin{align*}
1{,}331x&=161\,051+146\,410+133\,100=440\,561,\\
x&=\frac{440\,561}{1{,}331}=331\,000\text{ руб.}
\end{align*}
\labelp{Проверка} После первого года: $331\,000\cdot1{,}1-133\,100=231\,000$ руб. После второго: $231\,000\cdot1{,}1-133\,100=121\,000$ руб. После третьего: $121\,000\cdot1{,}1-133\,100=0$.
\result Первоначальный кредит --- 331\,000 руб.
\labelp{Задача на закрепление} Кредит погашен тремя равными ежегодными платежами по 66\,550 руб. при ставке $10\%$ годовых. Найдите первоначальную сумму кредита.

\task{9}{Переплата по кредиту}
\labelp{Текст задания} Для кредита из задания №\,7 найдите переплату в рублях и процент переплаты от первоначальной суммы. Процент округлите до сотых.
\labelp{Ваше решение} В работе записано вычитание $496{,}5-286{,}078\cdot3$ (тыс. руб.) с отрицательным результатом $-361{,}734$. Вычисление процента переплаты отсутствует.
\labelp{Ваш ответ} $-361\,734$ руб.; процент не указан.
\labelp{Где именно возникает ошибка} Последний корректный принцип --- сравнить общие платежи с первоначально взятой суммой. Первый неправильный шаг --- записать разность в обратном порядке и использовать ошибочный платёж из задания №\,7.
\err Переплата равна \emph{всему уплаченному} минус \emph{сумма полученного кредита}. Для данных задачи переплата должна быть положительной. Кроме того, условие требует процент от первоначальной суммы.
\idea Вначале найдите общий размер платежей, затем вычтите сумму кредита, после чего разделите переплату на первоначальный кредит и умножьте на $100\%$.
\labelp{Исправление} Используем корректный платёж из задания №\,7 --- 199\,650 руб.:
\[\Delta=3\cdot199\,650-496\,500.\]
\labelp{Полное правильное решение} Общая сумма выплат:
\[T=3\cdot199\,650=598\,950\text{ руб.}\]
Переплата:
\[\Delta=T-496\,500=102\,450\text{ руб.}\]
Процент переплаты:
\[q=\frac{102\,450}{496\,500}\cdot100\%\approx20{,}63\%.\]
\labelp{Проверка} Общая сумма выплат $598\,950$ руб. больше первоначального кредита $496\,500$ руб. на $102\,450$ руб.; процент рассчитан от исходной суммы кредита.
\result Переплата 102\,450 руб., или $20{,}63\%$.
\labelp{Задача на закрепление} Кредит 200\,000 руб. погашен тремя ежегодными платежами по 80\,000 руб. Найдите переплату в рублях и в процентах от первоначального кредита.

\task{10}{Четыре равных ежегодных платежа}
\labelp{Текст задания} Кредит погашен четырьмя равными ежегодными платежами по 146\,410 руб. Годовая ставка $10\%$. Найдите первоначальную сумму кредита, общий размер выплат и переплату.
\labelp{Ваше решение} В предоставленных пяти страницах решения задания №\,10 нет.
\labelp{Ваш ответ} Не указан.
\labelp{Где именно возникает ошибка} Решение не представлено; проверить собственные шаги невозможно. Главная недостающая часть --- уравнение, учитывающее все четыре платежа и ежегодный рост остатка долга.
\err В отсутствие вычислений задание остаётся невыполненным. Для восстановления метода важно не пропустить ни один платёж.
\idea Перенесём кредит и все выплаты к концу четвёртого года. Тогда первый платёж учитывается с коэффициентом $1{,}1^3$, второй --- $1{,}1^2$, третий --- $1{,}1$, четвёртый --- $1$.
\labelp{Исправление} Обозначим первоначальный кредит через $P$. При нулевом конечном долге:
\[P\cdot1{,}1^4=146\,410\,(1{,}1^3+1{,}1^2+1{,}1+1).\]
\labelp{Полное правильное решение} Вычислим коэффициенты:
\[1{,}1^4=1{,}4641,\qquad 1{,}331+1{,}21+1{,}1+1=4{,}641.\]
Поэтому
\[P=\frac{146\,410\cdot4{,}641}{1{,}4641}=464\,100\text{ руб.}\]
Общая сумма четырёх платежей:
\[T=4\cdot146\,410=585\,640\text{ руб.}\]
Переплата:
\[\Delta=T-P=585\,640-464\,100=121\,540\text{ руб.}\]
\labelp{Проверка} Остатки после платежей по годам: $464\,100\cdot1{,}1-146\,410=364\,100$; $364\,100\cdot1{,}1-146\,410=254\,100$; $254\,100\cdot1{,}1-146\,410=133\,100$; $133\,100\cdot1{,}1-146\,410=0$. Кредит погашен полностью.
\result Первоначальный кредит --- 464\,100 руб.; общая сумма выплат --- 585\,640 руб.; переплата --- 121\,540 руб.
\labelp{Задача на закрепление} Кредит выплачен четырьмя одинаковыми ежегодными платежами по 73\,205 руб. при годовой ставке $10\%$. Найдите первоначальную сумму кредита, все выплаты и переплату.
\end{document}
